\section{从 Language Model 到 AI Agent：一次调用为什么不够}

课堂在 49 分钟处切换到 agent。这个转场不是更换 buzzword，而是改变系统单位：language model 接受 token 并生成 token；agent 还要保持 state、读取 memory、选择 action、观察环境结果并根据反馈继续执行。一次模型调用可以是 agent 的 component，但不是完整 agent。

\subsection{Agent stack：model 外面还需要什么}

开场页列出 action、memory、communication 与 emergent architecture。Building AI Agents 页则把问题拆成 why 与 how：为什么需要 agent，怎样让模型进入现实接口。

\lecturefigure{slide-078.jpg}{AI Agent 部分的课程范围与研究方向}{官方 CS25 V4 slide deck，第 78 页}

\lecturefigure{slide-079.jpg}{Building AI Agents：why 与 how 的问题框架}{官方 CS25 V4 slide deck，第 79 页}

\begin{importantbox}{Agent 的最小闭环}
设环境状态为 $s_t$，模型根据 observation $o_t$、memory $m_t$ 与目标 $g$ 选择 action：
\[
a_t=\pi_\theta(o_t,m_t,g),\qquad
s_{t+1}=\mathcal{E}(s_t,a_t),\qquad
m_{t+1}=U(m_t,o_t,a_t,s_{t+1}).
\]
$\pi_\theta$ 是 policy，$\mathcal{E}$ 是环境转移，$U$ 是 memory update。Agent 必须处理行动结果，而不仅是预测下一个 token。
\end{importantbox}

\lecturefigure{slide-080.jpg}{单次 foundation-model call 缺少长期状态与行动反馈}{官方 CS25 V4 slide deck，第 80 页}

\lecturefigure{slide-081.jpg}{Agent architecture 在模型外增加 memory、tools 与 environment loop}{官方 CS25 V4 slide deck，第 81 页}

\begin{knowledgebox}{术语消化：Agent stack 的六个部件}
\begin{tabularx}{\textwidth}{L{0.19\textwidth} X}
\toprule
部件 & 职责 \\
\midrule
Model/policy & 解释 observation、提出计划或 action。 \\
State & 保存当前任务、页面、文件、变量和执行位置。 \\
Memory & 保存跨步骤或跨会话信息，并支持检索。 \\
Tool/action layer & 把结构化 action 映射到 API、browser、shell 或设备。 \\
Environment feedback & 返回成功、错误、页面变化、测试结果或传感器 observation。 \\
Verifier/guardrail & 检查约束、权限、结果与终止条件。 \\
\bottomrule
\end{tabularx}
\end{knowledgebox}

\teachervoice{讲者把转向 agent 的理由说得很直接：一次 foundation-model call 不够，因为真实任务还需要 memory、long context、personalization、actions、internet access 和跨步骤反馈。}

\subsection{Flight 与 driving-test demo：展示 action，不证明可靠性}

两页 demo 展示 MultiOn 早期 agent 预订航班和完成 California online driving test。它们说明模型可以把自然语言目标映射到浏览器操作，但公开视频中的单次成功不能推出一般成功率、支付安全或跨网站鲁棒性。读图时先辨认 observation、action 与完成条件，再区分演示证据和部署证据。

\lecturefigure{slide-082.jpg}{课堂展示的 AI agent 航班预订案例}{官方 CS25 V4 slide deck，第 82 页}

\lecturefigure{slide-083.jpg}{课堂展示的 online driving-test agent 案例}{官方 CS25 V4 slide deck，第 83 页}

\begin{warningbox}{Demo evidence 的正确读法}
Demo 可以证明“这条 action trajectory 至少发生过”，不能证明平均成功率、worst-case safety 或无人监督部署。需要补充任务集、重复次数、失败分类、权限范围、恢复策略与独立复现，才能从 capability demo 走向 engineering claim。
\end{warningbox}

\teachervoice{讲者把 driving-test demo 称为早期 real-world agent 展示。讲义保留这个历史事实，但不把它写成任何网站、账号或设备上都能稳定运行的结论。}

\subsection{为什么追求 human-like interface}

Human-like agent 的吸引力是复用现有界面：网页、desktop 和 mobile app 已为人设计。Agent 若能看、点击和输入，就不必等待每个服务提供 API；它也可能通过用户示范学习偏好。相同优势同时扩大风险，因为登录、支付和个人数据也暴露在 action space 中。

\lecturefigure{slide-084.jpg}{Human-like agent 复用人类界面、偏好与示范}{官方 CS25 V4 slide deck，第 84 页}

\begin{importantbox}{Interface coverage 与 safety boundary 是同一枚硬币}
更一般的 UI control 提高覆盖面，却削弱 typed API 提供的 schema、权限与 error code。设计时应优先使用受约束 API；只有在缺少接口时才扩大到 UI control，并为高风险 action 加 confirmation、least privilege、audit log 与 rollback/recovery。
\end{importantbox}

\subsection{Autonomy level 是责任梯度，不是营销等级}

课程借用 L0--L5 自动驾驶分级讨论 agent：低等级由人掌控，高等级逐步移除 human fallback。这个类比有助于说明责任变化，但不能直接用来给具体汽车或 agent 产品定级。下面的图要按“谁在监控、谁能接管、失败时谁负责”来读，而不是把等级当作单一能力分数。

\lecturefigure{slide-085.jpg}{从 human control 到 full autonomy 的五级类比}{官方 CS25 V4 slide deck，第 85 页}

\begin{warningbox}{课堂中的产品举例是非正式类比}
讲者口头把若干车辆系统对应到 L4/L5。正式 SAE 定义、运行设计域和当前产品状态比课堂类比严格得多，也会随时间变化。这里保留的教学结论只有一个：移除 human fallback 会显著提高验证与责任要求。
\end{warningbox}

\begin{knowledgebox}{Agent autonomy 的可操作分级}
\begin{tabularx}{\textwidth}{L{0.17\textwidth} X}
\toprule
等级 & 一个更适合软件 agent 的定义 \\
\midrule
Assist & 提供建议，不执行外部 action。 \\
Approve-each-step & 每个 action 都需用户确认。 \\
Bounded delegation & 在预定义 scope、预算和工具内自动执行。 \\
Supervised autonomy & 可长时间执行，但有实时 monitor、interrupt 与 fallback。 \\
Unsupervised autonomy & 无在线人类监督完成高影响任务；需要最严格的形式化边界和独立验证。 \\
\bottomrule
\end{tabularx}
\end{knowledgebox}

\subsection{API 与 direct computer interaction}

Agent 有两条主要 action 路线。API 提供结构化参数、返回值和权限；direct interaction 通过 screen、keyboard、mouse 或 accessibility tree 操作人类界面。后者覆盖面广，但更容易受 layout、latency、popup 和视觉歧义影响。

\lecturefigure{slide-087.jpg}{API control 与 direct computer interaction 两条路线}{官方 CS25 V4 slide deck，第 87 页}

\begin{importantbox}{Action schema}
一个可审计 action 不应只是自由文本，而应接近
\[
a_t=(\texttt{tool},\texttt{arguments},\texttt{precondition},\texttt{risk},\texttt{idempotency}).
\]
\texttt{precondition} 描述执行前状态，\texttt{risk} 决定是否需要确认，\texttt{idempotency} 说明重复调用是否安全。结构化 action 让 verifier 和 recovery 有明确对象。
\end{importantbox}

\lecturefigure{slide-088.jpg}{Universal Action API 与跨应用操作示例}{官方 CS25 V4 slide deck，第 88 页}

\begin{warningbox}{“不需要 API”不等于“API 没价值”}
UI agent 可以覆盖没有 API 的系统，但稳定 API 通常更快、更便宜、更容易授权和测试。理想架构是 capability routing：优先选择受约束工具，在必须使用 UI 时缩小可见区域、action set 与 credential scope。
\end{warningbox}

\subsection{本章小结}

Agent 把语言模型放进 observation--action--feedback loop，并增加 state、memory、tool、verifier 与 recovery。课堂 demo 展示了 action capability，autonomy 与 UI control 则暴露了责任和权限问题。真正的系统边界不在“模型会不会点击”，而在失败能否被发现、限制和恢复。

